\scanpage{143}
$H\to H^{\frac r2}$ de $L^{+\;\frac p2}$ dans $L^{\frac pr}$ ; or, les conditions de la prop. 1 sont aussitôt vérifiées \deleted{\unsure{en effet}} (condition 1 triviale ; condition 2 se voit \deleted{ainsi} en se ramenant à \deleted{\unsure{montrer que si}}\inserted{la dimension} finie, où c'est trivial ($\sum_{i=1}^{\unsure{n}}\lambda_i a_i\otimes a_i$, $\lambda_i\geq0$ et $a_i$ variables, $\to\lambda_i^{(0)}$ et $a_i^{(0)}$ ; on forme $\sum\lambda_i^{r/2}a_i\otimes a_i$ \unsure{\ldots{} --})

\uline{Remarque} 1 Si $0<p\leq2$, pour qu'une suite d'hermitiens $H_i$ dans $L^p$ converge \inserted{vers $H_0$}, il faut et il suffit qu'elle converge faiblement, et que $\|K_i\|\to\|H_0\|$ (ceci est connu au cas hilbertien)

% slanted
\margin{\deleted{\uline{\unsure{prendre comme|probablement}} \uline{\unsure{indépendant de $p$|pertinent ds $L^p$}} ! \unsure{la norme uniforme de Murray|On connaît le contre-exemple de Murray}}}

Est-ce vrai si $p<+\infty$ ? Uniforme convexité de $L^p$ ?

\uline{Remarque 2} Est-ce que $S_p\lvert H^r-K^r\rvert^{\frac1r}$ est fonction croissante de $r$ ($H$ hermitien positif) ? -- Il suffirait de savoir que $r<1$ implique

\[
\bigl\|\lvert H^r-K^r\rvert^{\frac1r}\bigr\|\leq\|H-K\|
\]

--- Alors on aurait une précision notable dans le corollaire 1.
% archive sheet number: 144
